% ===== PRÉAMBULE CANONIQUE — à coller VERBATIM en tête de CHAQUE .tex =====
\documentclass[11pt,a4paper]{article}
\usepackage[utf8]{inputenc}\usepackage[T1]{fontenc}\usepackage{lmodern}
\usepackage[french]{babel}
\usepackage{amsmath,amssymb,mathtools}\usepackage{icomma}
\usepackage[version=4]{mhchem}          % \ce{...} : équations & formules
\usepackage{siunitx}\sisetup{locale=FR} % \qty{}{}, \si{}, \ang{}
\usepackage{chemfig}                    % \chemfig{...} : structures développées
\usepackage{xcolor}\usepackage{colortbl}
\usepackage{tikz}\usepackage{pgfplots}\pgfplotsset{compat=1.18}
\usepackage[most]{tcolorbox}\usepackage{enumitem}\usepackage{array,booktabs}
\usepackage[a4paper,margin=2.2cm]{geometry}
\usetikzlibrary{arrows.meta,calc,angles,quotes,positioning,patterns,backgrounds,shapes.geometric}
\definecolor{primary}{RGB}{222,49,99}\definecolor{primarydark}{RGB}{150,22,55}
\definecolor{defcol}{RGB}{0,114,178}\definecolor{methcol}{RGB}{52,92,148}
\definecolor{excol}{RGB}{0,135,90}\definecolor{attncol}{RGB}{200,40,55}
\tcbset{boxbase/.style={enhanced,breakable,boxrule=0.9pt,arc=2.5pt,
  left=3mm,right=3mm,top=2.3mm,bottom=2.3mm,fonttitle=\bfseries,coltitle=white}}
% --- boîtes TUTORIEL (noms EXACTS, ne pas renommer) ---
\newtcolorbox{cle}[1][]{boxbase,colback=defcol!6,colframe=defcol,
  title={\textsc{À retenir}\ifx\relax#1\relax\relax\else\textmd{ : \itshape #1}\fi}}
\newtcolorbox{methode}[1][]{boxbase,colback=methcol!7,colframe=methcol,sharp corners,
  title={\textsc{Méthode}\ifx\relax#1\relax\relax\else\textmd{ --- \itshape #1}\fi}}
\newtcolorbox{exemplebox}[1][]{boxbase,colback=excol!5,colframe=excol,
  title={\textsc{Exemple}\ifx\relax#1\relax\relax\else\textmd{ : \itshape #1}\fi}}
\newtcolorbox{piege}[1][]{boxbase,colback=attncol!6,colframe=attncol,
  title={\textsc{Piège}\ifx\relax#1\relax\relax\else\textmd{ : \itshape #1}\fi}}
\newtcolorbox{remarque}[1][]{boxbase,colback=black!4,colframe=black!45,
  title={\textsc{Remarque}\ifx\relax#1\relax\relax\else\textmd{ : \itshape #1}\fi}}
% --- boîtes EXERCICE / CORRIGÉ (noms EXACTS, auto-comptées) ---
\newtcolorbox[auto counter]{exo}[1][]{boxbase,colback=primary!4,colframe=primary,
  title={\textsc{Exercice}~\thetcbcounter\ifx\relax#1\relax\relax\else\textmd{ --- \itshape #1}\fi}}
\newtcolorbox[auto counter]{corrige}[1][]{boxbase,colback=excol!4,colframe=excol,
  title={\textsc{Corrigé}~\thetcbcounter\ifx\relax#1\relax\relax\else\textmd{ --- \itshape #1}\fi}}
\newcommand{\R}{\mathbb{R}}\newcommand{\N}{\mathbb{N}}\newcommand{\Z}{\mathbb{Z}}\newcommand{\ds}{\displaystyle}
% (un \usepackage{fancyhdr}+en-tête est possible ; il est ignoré côté web)
% =========================================================================
\begin{document}
\sisetup{per-mode=symbol}
\begin{corrige}[chiffres non nuls]
Tout chiffre non nul est toujours significatif : on compte simplement les chiffres.
\begin{enumerate}[label=\alph*)]
  \item $837$ : $\mathbf{3}$ CS.
  \item $42,5$ : $\mathbf{3}$ CS (la virgule ne compte pas).
  \item $6,1234$ : $\mathbf{5}$ CS.
  \item $128,77$ : $\mathbf{5}$ CS.
\end{enumerate}
\end{corrige}

\begin{corrige}[zéros de tête]
Les zéros de tête (à gauche du premier chiffre non nul) ne servent qu'à caler la virgule : ils ne
sont \emph{jamais} significatifs. On compte à partir du premier chiffre non nul.
\begin{enumerate}[label=\alph*)]
  \item $0,00\mathbf{7}$ : $\mathbf{1}$ CS.
  \item $0,0\mathbf{52}$ : $\mathbf{2}$ CS.
  \item $0,00\mathbf{35}$ : $\mathbf{2}$ CS.
  \item $0,\mathbf{4}$ : $\mathbf{1}$ CS.
\end{enumerate}
\end{corrige}

\begin{corrige}[zéros captifs]
Un zéro captif (entre deux chiffres non nuls) est \emph{toujours} significatif.
\begin{enumerate}[label=\alph*)]
  \item $205$ : $\mathbf{3}$ CS.
  \item $4007$ : $\mathbf{4}$ CS (les deux zéros sont captifs).
  \item $3,05$ : $\mathbf{3}$ CS.
  \item $60,2$ : $\mathbf{3}$ CS.
\end{enumerate}
\end{corrige}

\begin{corrige}[zéros de queue après la virgule]
Un zéro de queue placé après la virgule est \emph{toujours} significatif : on ne l'écrirait pas s'il
n'était pas le fruit d'une mesure.
\begin{enumerate}[label=\alph*)]
  \item $15,200$ : $\mathbf{5}$ CS (les deux zéros finaux comptent).
  \item $0,0100$ : $\mathbf{3}$ CS (zéros de tête non sig.\ ; les deux zéros de queue le sont).
  \item $3,40$ : $\mathbf{3}$ CS.
  \item $8,000$ : $\mathbf{4}$ CS.
\end{enumerate}
\end{corrige}

\begin{corrige}[tous les cas à la fois]
On lit de gauche à droite en appliquant les trois règles sur les zéros.
\begin{enumerate}[label=\alph*)]
  \item $0,0035070$ : tête non sig.\ ; captif et queue sig.\ $\Rightarrow \mathbf{5}$ CS.
  \item $1030,0$ : la présence de la virgule rend \emph{tous} les chiffres significatifs
        (captif et zéro de queue) $\Rightarrow \mathbf{5}$ CS.
  \item $0,004080$ : tête non sig.\ ; le $0$ captif et le $0$ de queue sont sig.\ $\Rightarrow \mathbf{4}$ CS.
  \item $20,08$ : $\mathbf{4}$ CS (le $0$ captif et le $0$ après la virgule comptent).
\end{enumerate}
\end{corrige}

\end{document}
